\chapter{Using AI visibly}
\label{ch:ai-visibly}

\section*{What this chapter is for}

In a round where assistance is permitted, how you use it is assessed. This is
the one chapter in Part~\ref{part:build} that applies to only one of the three
formats, and it is the one where the assessment criteria are least familiar.

\section{What is actually scored}

Where criteria are published or reported, they converge on three things.

\textbf{Prompt quality, output validation, debugging.}
\reported[single]{At least one large employer's assisted round names exactly
these three.}{google-code-comprehension} Note what is absent: speed, and volume
of code produced.

\textbf{Catching the model's mistakes.} \reported[single]{Reporting on another
employer's assisted round suggests that finding and fixing what the assistant
got wrong is among the strongest positive
signals.}{meta-ai-enabled-coding}

\textbf{Iterative use rather than one large prompt.} \reported{Employers with
published policies make the same point \dash{} that dumping the whole problem
in and pasting the answer back is the behaviour they are watching for \dash{}
and interviewers are reported to be watching for reasoning rather than
copying.}{field-guide-trends}

\begin{kitnote}
Read those three together and the round is not measuring whether you can use
the tool. It is measuring whether you remain the engineer while using it. Every
one of the three is about your relationship to the output rather than about the
output.
\end{kitnote}

\section{Six habits}

\begin{kitmethod}
\textbf{1. Say what you are about to ask for, before you ask.} One sentence:
``I am going to ask it for the conflict-safe version of this update, because
that is the part I want to get right by hand rather than by memory.'' The
interviewer now knows what to watch for, and your intent is on the record
independently of what comes back.

\textbf{2. Ask for one thing.} A prompt that requests a feature produces
something you have to read for three minutes. A prompt that requests a function
produces something you can check in twenty seconds. The second is faster in
total and it is visibly under control.

\textbf{3. Read it out loud, at least once.} Not all of it, and not every time.
Once, on the part that matters, pointing at the screen. It is the clearest
possible demonstration of the validation criterion.

\textbf{4. Reject something, audibly.} This is the highest-value habit in the
chapter. When the output is wrong, or right but unsuitable, say why and
discard it. If nothing needs rejecting in a whole session, either you are
asking for things too small to be wrong, or you are not reading closely enough.

\textbf{5. Write the hard thing yourself.} The conflict case, the assertion,
the thing the design identified as the risk. Not out of pride \dash{} because
it is where your judgement is on display, and because it is the part you will
be asked about afterwards.

\textbf{6. Never paste what you cannot explain.} If it arrives and you do not
follow it, either work it out on screen or discard it. Accepting code you
cannot explain is the failure this round exists to detect, and it is detected
by one follow-up question.
\end{kitmethod}

\section{The Staff-level version}

Everything above is good practice at any level. One thing is specific to this
one.

\reported{Knowing when the model is the wrong tool is itself the signal, and a
borderline senior candidate is recorded as receiving a negative assessment for
repeatedly proposing model-based solutions to a problem that did not need
one.}{field-guide-trends}

\judgement{So the sentence with the most upside in the whole round is a refusal:
``I could do this part with the model and I am not going to, because it is
twenty lines of arithmetic that I want to be exactly right and I would rather
read it than review it.'' It demonstrates the judgement the level is being
assessed on, and it costs nothing.}

\section{Standards as a brake, not a plan}

If you work to conventions \dash{} your own or an employer's \dash{} bring them
as something the work is checked against, not as something to install.

A standard used as a brake is: writing a slice, then pausing to check it
against one rule, and saying which rule. A standard used as a plan is: spending
the first twenty minutes setting up the scaffolding it implies. The first is
judgement; the second is ceremony with a clock running, and it is the more
common failure among candidates who have prepared thoroughly.

\begin{kittrap}
The specific version of this failure: arriving with a project template and
spending ten minutes installing it before anything exists. The template is
usually a genuine asset and the ten minutes are usually fatal, because they are
the ten minutes in which discovery should have been turning into a running
skeleton.

Bring the template already set up if you can. If you cannot, bring a smaller
one.
\end{kittrap}

\section{When assistance is not permitted}

Then none of this applies, and one thing is worth saying: practise without it.
If your recent working habit is model-assisted, an unassisted round will be
slower than you expect in a way that is hard to predict from the inside, and
the rehearsal in Chapter~\ref{ch:rehearsal} should be run in whichever mode the
round will be.

\begin{drill}
Record yourself building something small with assistance, for thirty minutes.
Watch it back and count three things: prompts where you said in advance what
you wanted, outputs you read before accepting, and outputs you rejected. If the
third count is zero, that is the habit to work on, and it is the one being
scored most heavily.
\end{drill}

\section*{What this chapter does not claim}

Two of the three published criteria rest on single sources and are marked as
such in the text. The third is better supported, coming from several employers'
stated policies.

The six habits are the author's, assembled from what those criteria imply. They
have not been measured against outcomes, and a round that scores differently
from the published descriptions would reward different behaviour.
